\documentclass{eeleyes}
\usepackage{subfiles}

\usepackage{fancyhdr}
\usepackage{dsfont}

\newcommand\conj[1]{\overline{#1}}
\newcommand\term[1]{\textbf{#1}}
\newcommand\eps{\varepsilon}
\DeclareMathOperator{\interior}{int}

\newcommand\thh{\text{th}}
\newcommand{\coloneqcolon}{\mathrel{\mathop{:}\!\!=\!\!\mathop{:}}}

\NewDocumentEnvironment{romanlist}{}
{\begin{enumerate}[label=\roman*)]}
{\end{enumerate}}

\begin{document}

% \tableofdefs
% \newpage

\newtheorem{com}{{\color{blue} Comment}}[section]
\newtheorem{que}{{\color{blue} Question}}[section]
\newtheorem*{ans}{{\color{orange}\em Answer}}
\newtheorem*{rem}{{\color{magenta}{\em Remark}}}
\newtheorem*{hint}{{\color{brown}{\em Hint}}}

\pagestyle{fancy}

\subfile{dis-sep-30.tex}

\begin{rem}
    \newcommand{\ind}[1]{\mathds{1}_{#1}}
    Different non-decreasing functions can generate the same measure. Consider $f = \ind{[0,\infty)}$ and $g = \ind{(0,\infty)}$. Then $f(0) \neq g(0)$. Take $a < b$ and consider the following cases
    \begin{itemize}
        \item[$a,b \leq 0)$]
            Both $f(a) = 0$ and $g(b) = 0$, thus $\mu_f((a,b]) = \mu_g((a,b]) = 0$.

        \item[$a,b \geq 0)$]
            Since $0 \leq a < b$, then $f(a) = f(b) = 1$ and $g(a) = g(b) = 1$, thus $\mu_f((a,b]) = \mu_g((a,b]) = 0$.

        \item[$a < 0 < b)$]
            Note that
            \begin{alignat*}{5}
                \mu_f((a,b]) &= f(b) - f(a) &&= 1 - 0 &&= 1 \\
                \mu_g((a,b]) &= g(b) - g(a) &&= 1 - 0 &&= 1
            \end{alignat*}
    \end{itemize}
    Therefore $\mu_f = \mu_g$ despite $f \neq g$. 
\end{rem}

\begin{theorem}
    If $\mu$ is a Borel measure that is finite for all bounded intervals, then there is a unique (up to an additive constant) non-decreasing right-continuous function $F : \R \to \R$ such that
    \[
        \mu((a,b]) = F(b) - F(a), \quad a < b
    .\]
\end{theorem}

\begin{proof}
    Define the candidate generating function
    \[
        F(x) \coloneq \begin{cases}
            \mu((0,x]) & x \geq 0 \\
            -\mu((x, 0]) & x < 0 \\
        \end{cases}
    .\]
    Since $\mu((x,y]) < \infty$ for any $x,y \in \R$, it follows that $F(x) < \infty$ for all $x$. Suppose then $a < b$ and consider the following cases
    \begin{itemize}
        \item[$a,b \geq 0$)]
            Since $(a,b] = (0, b] \setminus (0,a]$, then
            \[
                F(b) - F(a) = \mu((0,b]) - \mu((0,a]) = \mu((a,b])
            .\]

        \item[$a,b < 0$)]
            Since $(a,b] = (a,0] \setminus (b, 0]$, then
            \begin{align*}
                F(b) - F(a) &= (-\mu((b,0])) - (-\mu((a,0])) \\
                &= \mu((a,0]) - \mu((b, 0]) \\
                &= \mu((a,b])
            \end{align*}

        \item[$a < 0 \leq b$)]
            Since $(a,b] = (a,0] \;\dot{\cup}\; (0,b]$, then
            \begin{align*}
                F(b) - F(a) &= \mu((0,b]) - (-\mu((a,0])) \\
                &= \mu((0,b]) + \mu((a,0]) \\
                &= \mu((a,b])
            \end{align*}
    \end{itemize}
    Clearly $F$ is non-decreasing since $F(y) - F(x) = \mu((x,y]) \geq 0$ for $y > x$. Take $x \in \R$ and $(h_n)_{n \in \N}$ with $h_n \geq 0$ and $\lim_{n \to \infty} h_n = 0$. Note then that
    \[
        F(x + h_n) - F(x) = \mu((x, x+h_n])
    \]
    which by continuity from above means $\lim_{n \to \infty} \mu((x,x+h_n]) = 0$ and hence $\lim_{n \to \infty} F(x + h_n) = F(x)$. Thus $F$ is right-continuous.

    For uniqueness, suppose there is some $G$ that is also non-decreasing and right-continuous such that for $a < b$
    \[
        \mu((a,b]) = G(b) - G(a)
    .\]
    Then for all $a < b$
    \[
        F(b) - F(a) = G(b) - G(a) \implies F(b) - G(b) = F(a) - G(a)
    .\]
    This means that $F - G$ is a constant function, hence $G = F + c$ for some constant $c$.
\end{proof}

\end{document}

